\documentclass[a4paper]{article}
\usepackage[pdftex]{graphicx}
\usepackage[utf8]{inputenc}
\usepackage{enumerate}
\usepackage{icomma}
\usepackage{amssymb}
\usepackage{tikz}
\usepackage{href-ul}
\hypersetup{
	colorlinks=true,
	linkcolor=blue,
	urlcolor=blue}
\usepackage{geometry}
\geometry{a4paper, top=15mm, left=15mm, right=15mm, bottom=15mm,
	headsep=10mm, footskip=12mm}

\begin{document}
\begin{tikzpicture}(10,3)
	\draw[ultra thick](2,0) --(10,0) -- (10,1.3) --(2,1.3) -- (2,0);
	\draw[fill=black](2,0)-- (0,.6) -- (2,.6) -- (2,0);
	\node at (6,.6) {\large Lösungsblatt von \href{https://www.okuyakl.de}{www.okuyakl.de}};
\end{tikzpicture}
\vspace{0.5 cm}

\noindent{\bf Aufgabe 1.}\\
$$\overrightarrow{AB} =\left( \begin{array}{c} 2,5 - (-2) \\ -3-3 \end{array}\right)= \left( \begin{array}{c} 4,5 \\ -6 \end{array}\right)$$
$$\overrightarrow{AC} =\left( \begin{array}{c} 6 - (-2) \\ 7,5-3 \end{array}\right)= \left( \begin{array}{c} 8 \\ 4,5 \end{array}\right)$$
Den Flächeninhalt berechnen wir mit der Determinante:
$$A = {1 \over 2}|\overrightarrow{AB} \times \overrightarrow{AB}| ={1 \over 2} \left| \begin{array}{cc} 4,5 & 8 \\ -6 & 4,5 \end{array}\right|= {1 \over 2} \cdot | 20,25 + 48 | = 34,125~{\rm FE} $$

\noindent {\bf Aufgabe 2. a)}

Zunächst berechnen wir $\overrightarrow{AC}$ und $\overrightarrow{AD}$:

$$
\renewcommand{\arraystretch}{2}
\begin{array}{rcl}
 \overrightarrow{AC}&={ 10-2 \choose 7-3} &= {8 \choose 4} \\
 \overrightarrow{AD}&={  x-2 \choose -0,5x+5-3} &= {x-2 \choose -0,5x+2}
\end{array}
$$

Dann stellen wir die Determinante auf:

$$ 
\renewcommand{\arraystretch}{2}
\begin{array}{rcl}
 A(x)= \left| \begin{array}{cc}  8 & x-2 \\ 4 & -0,5x+2 \end{array} \right|
&=& 8 \cdot (-0,5x+2) - 4(x-2)\\
A(x)&=&-4x +16 -4x + 8 = (-8x+24) ~{\rm FE}
\end{array}
$$

\noindent {\bf Aufgabe 2. b)}

Wir setzen gleich:

$$ 
\begin{array}{rcl}
A(x)= -8x+24&=& 16\\
    x&=& 1 
 \end{array}
$$
Damit sind die Koordinaten $D(1|4,5)$

\noindent {\bf Aufgabe 3.}

Bestimmung von $\overrightarrow{AB}$  und  $\overrightarrow{AC}$:

$$
\renewcommand{\arraystretch}{2}
\begin{array}{rcl}
 \overrightarrow{AB}&={ x_B-2 \choose 4-2} &= {x_B -2 \choose 2} \\
 \overrightarrow{AC}&={ 10-2 \choose 8-2} &= {8 \choose 6}
\end{array}
$$

Bestimmung von $x_B$ mittels der Determinante:

$$ 
\renewcommand{\arraystretch}{2}
\begin{array}{rcl}
 {1 \over 2} \left| \begin{array}{cc}  x_B-2&8 \\ 2&6 \end{array} \right|&=& {1 \over 2} \cdot (6( x_B-2) - 16) \\
3 x_B - 6 - 8 &=& 19 \\
3 x_B &=& 33 \\
x_B &=& 11
\end{array}
$$

\fbox{
	\begin{minipage}{0.4\textwidth}
		\noindent{\bf Aufgabe 4. a)}\\
		\includegraphics[width=7 cm]{drck256}	
	\end{minipage}
}
\hfill
\fbox{
	\begin{minipage}{0.55\textwidth}
		\noindent{\bf Aufgabe 4. b)}\\
		Fläche = Grundseite mal Höhe :2, Grundseite = 2x-(-x+3);\\
		Höhe = 6-x.
		$$A= {1 \over 2} \cdot (3x-3) \cdot (6-x) $$
		$$= {1\over 2} (-3x^2 + 18x +3x -18) = -1,5x^2 + 10,5x -9 \quad {\rm FE} $$
		
		\noindent{\bf Aufgabe 4. c)}\\
		Quadratische Ergänzung:
		$$
		\renewcommand{\arraystretch}{2}
		\begin{array}{rcl}
		A(x) &=& -1,5(x^2 - 7x + 3,5^2-3,5^2)-9 \\ 
		&=& -1,5[(x-3,5)^2 -12,25]-9 \\
		&=& -1,5(x-3,5)^2 +9,375 \\
		\hline
		A_{max}&=& 9,375 \quad \textnormal{für} \quad x=3,5
		\end{array}
		$$	
	\end{minipage}
}

\noindent {\bf Aufgabe 5. a)}
Der Flächeninhalt eines Rechteckes ist $A=\overline{OB_n} \cdot \overline{B_nC_n}$
Mit $\overline{OB_n}=x$ und $\overline{B_nC_n}=-0,4x+4,9$ erhalten wir:
$$A = x \cdot (-0,4x+4,9)= -0,4x^2 +4,9x$$

\noindent {\bf Aufgabe 5. b)}
Extremwertbestimmung mittels quadratischer Ergänzung liefert:
$$A_{max}= 15,01 \qquad \textnormal{für} \qquad x=6,125$$ 

\noindent {\bf Aufgabe 5. c)}
Für den Eckpunkt $C_2$  des Quadrates $OB_2C_2D_2$ gilt: x- und y-Koordinate sind gleich. Wir lösen:
$$
\renewcommand{\arraystretch}{2}
\begin{array}{rcll}

x&=&-0,4x+ 4,9 & |+ 0,4x  \\
1,4x &=& 4,9 & |:1,4 \\
x  &=& 3,5 &\Rightarrow \quad C_2(3,5|3,5)
\end{array}
$$

\noindent {\bf Aufgabe 5. d)}
Die Seitenlänge des neuen Quadrates ist $x-3,5$
Für den Eckpunkt $C_3$ des Quadrates $B_2EC_3F$ gilt:
$$
\renewcommand{\arraystretch}{2}
\begin{array}{rcll}

 x-3,5 &=& -0,4x +4,9  \\
 1,4 x &=& 8,4 \\
 x &=& 6  & \Rightarrow \quad C_3(6|2,5)
\end{array}
$$
Das neue Quadrat hat also die Seitenlänge 2,5.

\includegraphics[width=6 cm]{quger038}

\begin{center}
	\includegraphics[width=7 cm]{../../viecher/endcomic.pdf}
	
	Hier geht es zurück zum \href{https://www.okuyakl.de/math/m9kodiL038/aa038.pdf}{Aufgabenblatt}
\end{center}

\end{document}

